\chapter{Тестирование}
\label{ch:testing}

\section{Развернутое описание тестирования проекта онлайн-кинотеатра}

Тестирование проекта онлайн-кинотеатра было комплексным и включало в себя множество этапов, начиная от модульного тестирования отдельных компонентов бэкенда и заканчивая нагрузочным тестированием всей системы. Основной целью тестирования было обеспечение стабильности, производительности и корректности работы всех компонентов системы, включая API, базу данных и ML-модель для рекомендаций. Ниже приведено детальное описание каждого этапа тестирования, использованных метрик, проблем, с которыми столкнулись, и изменений, внесенных после тестирования.

\section{Тестирование кода}

\subsection*{Модульное тестирование (Unit Testing)}
Модульное тестирование было первым этапом проверки корректности работы кода. Для этого использовались фреймворки, такие как pytest (для Python) или JUnit (для Java, если применялся Spring). Каждый модуль, отвечающий за CRUD-операции (создание, чтение, обновление, удаление данных о фильмах, тегах и пользователях), был протестирован на корректность выполнения операций.

\begin{itemize}
    \item Пример тестов:
    \begin{itemize}
        \item Создание фильма: Проверялось, что фильм с корректными данными (название, описание, теги) успешно создается и сохраняется в базе данных. Также проверялась обработка ошибок при попытке создания фильма с некорректными данными (например, отсутствие названия или неверный формат данных).
        \item Чтение данных: Проверялось, что данные о фильме корректно извлекаются из базы данных, включая связанные теги и информацию о рекомендациях.
        \item Обновление данных: Проверялось, что данные о фильме могут быть обновлены, включая изменение названия, описания или тегов. Также проверялась обработка ошибок при попытке обновления несуществующего фильма.
        \item Удаление данных: Проверялось, что фильм и связанные с ним теги успешно удаляются из базы данных. Также проверялась обработка ошибок при попытке удаления несуществующего фильма.
    \end{itemize}
\end{itemize}

\subsection*{Интеграционное тестирование (Integration Testing)}
Интеграционные тесты проверяли взаимодействие между различными компонентами системы, такими как база данных, API и ML-модель. Эти тесты были направлены на проверку корректности работы системы в целом.

\begin{itemize}
    \item Пример тестов:
    \begin{itemize}
        \item Загрузка рекомендаций: Проверялось, что данные о рекомендациях, сгенерированные ML-моделью, корректно загружаются в базу данных и могут быть получены через API.
        \item Обновление данных: Проверялось, что обновление данных о фильме (например, изменение тегов) корректно отражается в рекомендациях.
        \item Взаимодействие с API: Проверялось, что API корректно обрабатывает запросы на получение данных о фильмах, тегах и рекомендациях, а также возвращает корректные ошибки при некорректных запросах.
    \end{itemize}
\end{itemize}

\subsection*{Нагрузочное тестирование (Load Testing)}
Для проверки устойчивости системы к высокой нагрузке использовались инструменты, такие как Apache JMeter или Locust. Эти тесты были направлены на проверку производительности системы при большом количестве одновременных запросов.

\begin{itemize}
    \item Пример тестов:
    \begin{itemize}
        \item Получение данных о фильмах: Проверялось время отклика API при одновременном выполнении 1000 запросов на получение данных о фильмах.
        \item Генерация рекомендаций: Проверялось время отклика при генерации рекомендаций для большого количества пользователей (например, 10 000 пользователей).
        \item Обновление данных: Проверялось время выполнения запросов на обновление данных о фильмах при высокой нагрузке.
    \end{itemize}
\end{itemize}

\subsection*{Тестирование безопасности (Security Testing)}
Проверялась защищенность API от распространенных уязвимостей, таких как SQL-инъекции, XSS (межсайтовый скриптинг) и CSRF (межсайтовая подделка запроса).

\begin{itemize}
    \item Пример тестов:
    \begin{itemize}
        \item SQL-инъекции: Проверялось, что система корректно обрабатывает попытки внедрения SQL-кода через входные данные (например, поле названия фильма).
        \item XSS: Проверялось, что система корректно обрабатывает попытки внедрения вредоносного JavaScript-кода через входные данные.
        \item CSRF: Проверялось, что система корректно обрабатывает попытки выполнения несанкционированных действий от имени пользователя.
        \item Авторизация и аутентификация: Проверялось, что только авторизованные пользователи могут выполнять определенные действия (например, обновление данных о фильмах).
    \end{itemize}
\end{itemize}

\section{Тестирование ML-модели}

\subsection*{Тестирование качества рекомендаций}
ML-модель, генерирующая рекомендации, тестировалась на основе исторических данных о просмотрах и предпочтениях пользователей. Основной целью было оценить, насколько рекомендации соответствуют предпочтениям пользователей.

\begin{itemize}
    \item Метрики:
    \begin{itemize}
        \item Precision@K и Recall@K — для оценки точности и полноты рекомендаций. Precision@K показывает, какая доля рекомендованных фильмов действительно интересна пользователю, а Recall@K — какая доля интересных фильмов была рекомендована.
        \item RMSE (Root Mean Squared Error) — для оценки ошибки предсказания рейтингов. Эта метрика использовалась для оценки точности предсказания рейтингов фильмов.
        \item MAP (Mean Average Precision) — для оценки средней точности рекомендаций. Эта метрика учитывает порядок рекомендаций и показывает, насколько хорошо модель ранжирует фильмы.
        \item Coverage — для оценки доли фильмов, которые модель может рекомендовать. Эта метрика показывает, насколько широкий охват имеет модель.
    \end{itemize}
\end{itemize}

\subsection*{A/B тестирование}
Для оценки эффективности рекомендаций проводилось A/B тестирование, где одна группа пользователей получала рекомендации от старой модели, а другая — от новой.

\begin{itemize}
    \item Метрики:
    \begin{itemize}
        \item CTR (Click-Through Rate) — процент пользователей, которые кликнули на рекомендованный фильм. Эта метрика показывала, насколько рекомендации привлекательны для пользователей.
        \item Среднее время просмотра фильма: Эта метрика показывала, насколько рекомендованные фильмы соответствуют интересам пользователей.
    \end{itemize}
\end{itemize}

\subsection*{Проблемы при тестировании ML-модели}
\begin{itemize}
    \item Холодный старт: Проблема с рекомендациями для новых пользователей или новых фильмов, по которым недостаточно данных.
    \item Смещение данных: Исторические данные могли быть смещены в сторону популярных фильмов, что снижало качество рекомендаций для нишевого контента.
    \item Низкая точность рекомендаций: В некоторых случаях модель рекомендовала фильмы, которые не соответствовали интересам пользователей.
\end{itemize}

\subsection*{Изменения после тестирования}
\begin{itemize}
    \item Добавление гибридной модели, сочетающей коллаборативную фильтрацию и контент-базированные методы. Это позволило улучшить качество рекомендаций для нишевого контента.
    \item Внедрение механизма обработки холодного старта (например, рекомендации на основе популярных фильмов или тегов).
    \item Оптимизация алгоритма обучения модели для повышения точности рекомендаций.
\end{itemize}

\section{Тестирование API}

Тестирование API было одним из ключевых этапов проверки корректности работы системы онлайн-кинотеатра. API выступает в роли основного интерфейса взаимодействия между клиентской частью (например, веб-приложением или мобильным приложением) и серверной частью, включая базу данных и ML-модель. Основной целью тестирования API было обеспечение стабильности, производительности и безопасности всех эндпоинтов. Ниже приведено детальное описание

\section{Функциональное тестирование API}

Функциональное тестирование было направлено на проверку корректности работы всех эндпоинтов API. Каждый эндпоинт тестировался на соответствие ожидаемому поведению, включая обработку корректных и некорректных запросов.

\subsection*{Тестирование эндпоинтов для работы с фильмами}
API предоставляет несколько эндпоинтов для работы с фильмами, включая создание, чтение, обновление и удаление данных.

\begin{itemize}
    \item Создание фильма:
    \begin{itemize}
        \item Проверялось, что фильм с корректными данными (название, описание, теги) успешно создается и сохраняется в базе данных.
        \item Проверялась обработка ошибок при попытке создания фильма с некорректными данными (например, отсутствие названия или неверный формат данных).
        \item Проверялась обработка дубликатов (например, попытка создания фильма с уже существующим названием).
    \end{itemize}
    \item Чтение данных о фильме:
    \begin{itemize}
        \item Проверялось, что данные о фильме корректно извлекаются из базы данных, включая связанные теги и информацию о рекомендациях.
        \item Проверялась обработка ошибок при попытке получения данных о несуществующем фильме.
        \item Проверялась корректность пагинации и фильтрации (например, получение списка фильмов по определенному тегу).
    \end{itemize}
    \item Обновление данных о фильме:
    \begin{itemize}
        \item Проверялось, что данные о фильме могут быть обновлены, включая изменение названия, описания или тегов.
        \item Проверялась обработка ошибок при попытке обновления несуществующего фильма.
        \item Проверялась корректность обновления связанных данных (например, тегов).
    \end{itemize}
    \item Удаление фильма:
    \begin{itemize}
        \item Проверялось, что фильм и связанные с ним теги успешно удаляются из базы данных.
        \item Проверялась обработка ошибок при попытке удаления несуществующего фильма.
    \end{itemize}
\end{itemize}

\subsection*{Тестирование эндпоинтов для работы с тегами}
API также предоставляет эндпоинты для работы с тегами, которые используются для категоризации фильмов.

\begin{itemize}
    \item Создание тега:
    \begin{itemize}
        \item Проверялось, что тег с корректными данными (название) успешно создается и сохраняется в базе данных.
        \item Проверялась обработка ошибок при попытке создания тега с некорректными данными (например, отсутствие названия).
    \end{itemize}
    \item Чтение данных о теге:
    \begin{itemize}
        \item Проверялось, что данные о теге корректно извлекаются из базы данных.
        \item Проверялась обработка ошибок при попытке получения данных о несуществующем теге.
    \end{itemize}
    \item Обновление данных о теге:
    \begin{itemize}
        \item Проверялось, что данные о теге могут быть обновлены, включая изменение названия.
        \item Проверялась обработка ошибок при попытке обновления несуществующего тега.
    \end{itemize}
    \item Удаление тега:
    \begin{itemize}
        \item Проверялось, что тег успешно удаляется из базы данных.
        \item Проверялась обработка ошибок при попытке удаления несуществующего тега.
    \end{itemize}
\end{itemize}

\subsection*{Тестирование эндпоинтов для работы с рекомендациями}
API предоставляет эндпоинты для получения рекомендаций для пользователей и фильмов.

\begin{itemize}
    \item Получение рекомендаций для пользователя:
    \begin{itemize}
        \item Проверялось, что API корректно возвращает список рекомендованных фильмов для конкретного пользователя.
        \item Проверялась обработка ошибок при попытке получения рекомендаций для несуществующего пользователя.
    \end{itemize}
    \item Получение рекомендаций для фильма:
    \begin{itemize}
        \item Проверялось, что API корректно возвращает список фильмов, рекомендованных на основе выбранного фильма.
        \item Проверялась обработка ошибок при попытке получения рекомендаций для несуществующего фильма.
    \end{itemize}
\end{itemize}


\endinput